\chapter{The Administrator Portal}
\label{ch:admin}

\portalcover{Administrator Portal}{For plant management and the system administrator}{secDisp}

\section{Who this chapter is for}

You either run the plant or look after the system. You can see everything and
change everything, which is why this chapter also explains what \emph{not} to do.

\begin{table}[H]
\centering
\small
\begin{tabular}{@{}p{40mm}p{98mm}@{}}
\toprule
\textbf{Your job} & \textbf{The screens you will live in} \\
\midrule
Plant Manager &
  Dashboard, Section WIP, Orders, Reports \\
System Administrator &
  Users \& Readers, Master Data, Audit Trail \\
\bottomrule
\end{tabular}
\end{table}

\section{The dashboard}

\menu{Dashboard} is the whole factory on one page. It refreshes every time you
open it.

\shot{10-admin-dashboard}{The plant dashboard.}{fig:adm-dash}

\subsection{The eight headline numbers}

\begin{table}[H]
\centering
\small
\begin{tabular}{@{}p{34mm}p{104mm}@{}}
\toprule
\textbf{Number} & \textbf{What it really tells you} \\
\midrule
Work in progress   & Every garment on the floor, not yet shipped. \\
Tagged today       & New garments registered at stitching --- your real input rate. \\
Shipped today      & Units that have left the building. \\
In transit         & Garments between two departments right now. A high number means trolleys are sitting around. \\
QC pass rate       & Percentage passed today. \\
Open rework        & Garments waiting to be corrected. \\
Variances          & Handovers where the count did not agree. \textbf{This should be zero.} \\
Fabric in store    & Metres available and how many rolls. \\
\bottomrule
\end{tabular}
\end{table}

\subsection{The section cards}

Each card is a department. The big number is what it is holding; the coloured bar
underneath is the ageing profile. Click any card to drill into it.

\begin{plainwords}
The fabric warehouse and cutting are measured in \emph{rolls} and \emph{bundles},
not garments, because garments do not exist yet at that point in the process.
That is why their cards show different units.
\end{plainwords}

\subsection{Output over the last 24 hours}

The bar chart shows garments tagged, dispatched and QC-passed for each hour, and
the table under it compares the last five days shift by shift. This is where you
see the difference between the morning and evening shift.

\subsection{What needs attention}

Two panels list problems: handovers that did not tally, and work that has been
sitting more than 24 hours. Both link straight to the detail.

\section{Orders and how production is progressing}

Open \menu{Orders}.

\shot{14-admin-orders}{Customer orders with live progress.}{fig:adm-orders}

Click any order to see it size by size.

\shotsmall{15-admin-order-detail}{One order, broken down by size, showing how many
are tagged, how many shipped, and where the rest currently are.}{fig:adm-order}

The \field{Where it is now} column is the useful one: it shows how many pieces of
that size are sitting in each department at this moment.

\section{Reports}

The report builder is described step by step in Chapter~\ref{ch:reports}. As
an administrator you also have these datasets:

\begin{table}[H]
\centering
\small
\begin{tabular}{@{}p{40mm}p{98mm}@{}}
\toprule
\textbf{Dataset} & \textbf{One row is\ldots} \\
\midrule
Articles / WIP        & one garment \\
Tracking Events       & one scan or status change \\
Dispatch \& Receipt   & one handover document \\
QC Inspections        & one inspection (or one defect, if a defect column is used) \\
Fabric Rolls          & one roll \\
Cutting \& Bundles    & one bundle \\
Shipments             & one packed garment \\
\bottomrule
\end{tabular}
\end{table}

\begin{tipbox}
Press \btn{Show query} beside any result to see exactly what the system asked the
database. Useful when someone questions a number.
\end{tipbox}

\section{Master data}

\menu{Master Data} holds the reference lists everything else depends on.

\shot{16-admin-masters-styles}{Designs, with the picture QC marks faults on.}{fig:adm-masters}

\begin{table}[H]
\centering
\small
\begin{tabular}{@{}p{34mm}p{104mm}@{}}
\toprule
\textbf{Tab} & \textbf{What to keep correct} \\
\midrule
Designs / Styles &
  The front and back images. \textbf{These are what QC clicks on}, so an
  out-of-date picture means faults get marked in the wrong place. \\
Customers &
  The tag specification. This is shown to the dispatch operator when they apply
  the customer's tag. \\
Fabric types, Colours &
  Keep codes stable. Renaming a colour changes it everywhere, including history. \\
Sizes &
  The sort order controls how sizes appear on every report and document. Set it
  once, properly (28, 30, 32\ldots\ not alphabetically). \\
Defect codes &
  The list QC chooses from. Too many codes and inspectors pick at random; too
  few and you learn nothing. \\
\bottomrule
\end{tabular}
\end{table}

\shotsmall{17-admin-masters-edit}{Editing a design.}{fig:adm-master-edit}

\begin{warnbox}
Set items to \emph{Inactive} rather than deleting them. Inactive items stop
appearing in dropdowns but stay attached to the history that already refers to
them.
\end{warnbox}

\section{User accounts}

Open \menu{Users \& Readers}.

\shot{18-admin-users}{Staff accounts.}{fig:adm-users}

\begin{steps}
  \item Press \btn{+ New user}.
  \item Enter the username, full name and employee number.
  \item Choose the \field{Role} --- this decides what they can do.
  \item Choose the \field{Home section} --- this decides which transfers they see
        first when they sign in.
  \item Set an initial password and have them change it.
\end{steps}

\subsection{Choosing the right role}

\shot{20-admin-roles}{The permission matrix, published by the server itself, so it
always matches what is actually enforced.}{fig:adm-roles}

\begin{table}[H]
\centering
\small
\begin{tabular}{@{}p{30mm}p{108mm}@{}}
\toprule
\textbf{Role} & \textbf{Give it to} \\
\midrule
\texttt{OPERATOR}   & Station staff. Can scan, tag, sort, receive and dispatch. \\
\texttt{QC}         & Inspectors. Can inspect and move batches, nothing else. \\
\texttt{SUPERVISOR} & Department heads. Adds closing variances and cancelling dispatches. \\
\texttt{MANAGER}    & Plant management. All the numbers and reports, plus orders and master data --- but not day-to-day scanning. \\
\texttt{ADMIN}      & You and your deputy only. \\
\texttt{VIEWER}     & Auditors and visitors. Can look, cannot touch. \\
\bottomrule
\end{tabular}
\end{table}

\begin{stopbox}
\textbf{Give the smallest role that lets someone do their job.} Every action is
recorded against the person who did it, which is only meaningful if people sign
in as themselves and cannot do things outside their job.
\end{stopbox}

\subsection{Disabling somebody who has left}

Set their status to \emph{Disabled}. Do not delete the account --- their name must
remain attached to the work they did.

\section{RFID readers}

The \menu{RFID readers} tab is where fixed readers are registered.

\shot{19-admin-readers}{Registered readers, and the instructions for connecting
one.}{fig:adm-readers}

\begin{steps}
  \item Press \btn{+ New reader}. Give it a code, a name, the section it belongs
        to and its type.
  \item Press \btn{Issue key}. A long key is shown \textbf{once}.
  \item Configure the reader to send that key with every read.
\end{steps}

\begin{table}[H]
\centering
\small
\begin{tabular}{@{}p{28mm}p{110mm}@{}}
\toprule
\textbf{Type} & \textbf{Where it goes} \\
\midrule
\texttt{HANDHELD} & Carried by an operator. Needs no key. \\
\texttt{PORTAL}   & A doorway. Reads what passes through. \\
\texttt{TUNNEL}   & A conveyor. Best for reading dense piles. \\
\texttt{TABLETOP} & A bench, for one garment at a time. \\
\texttt{ENCODER}  & Writes new tags --- stitching and dispatch. \\
\bottomrule
\end{tabular}
\end{table}

\begin{warnbox}
The key is shown once and never again. If it is lost, press \btn{Reissue key} ---
the old key stops working immediately, so update the reader at the same time.
\end{warnbox}

The \field{Last seen} column tells you whether a reader is actually talking to the
system. A reader that has not been seen since yesterday needs a look.

\section{The audit trail}

\menu{Audit Trail} records every change, permanently.

\shot{21-admin-audit}{The audit trail, filtered and searchable.}{fig:adm-audit}

You can filter by user, action, record type and date. Click any entry for the
full record.

\begin{table}[H]
\centering
\small
\begin{tabular}{@{}p{36mm}p{102mm}@{}}
\toprule
\textbf{Look for} & \textbf{When} \\
\midrule
\texttt{LOGIN\_FAILED}   & Somebody cannot get in, or somebody is guessing passwords. \\
\texttt{VARIANCE\_CLOSE} & A supervisor wrote off missing pieces. Read the reason. \\
\texttt{ARTICLE\_ADJUST} & A garment was moved by hand. \\
\texttt{TAG\_SWAP}       & A tag was replaced outside dispatch. \\
\texttt{SCRAP}           & A garment was written off. \\
\bottomrule
\end{tabular}
\end{table}

\section{Tracing any garment}

\menu{Trace Article} answers ``what happened to this one?''

\shot{98-trace-search}{Search by tag or serial number.}{fig:adm-trace-search}

Type or scan any tag the garment has \emph{ever} carried, or its serial number.
The result is the complete life story: every scan, in order, with the person, the
shift, the reader and the document.

\begin{plainwords}
If you scan a tracking tag that was removed at dispatch, the system says ``not
registered''. That is correct, not a fault. That tag has been released and may
already be on a different garment, so it would be wrong to claim it still
identifies the old one. Search by serial number instead.
\end{plainwords}

\section{Routine housekeeping}

\begin{table}[H]
\centering
\small
\begin{tabular}{@{}p{26mm}p{112mm}@{}}
\toprule
\textbf{How often} & \textbf{What to do} \\
\midrule
Daily   & Check \emph{Variances} on the dashboard is zero. Chase any that are not. \\
Daily   & Check the \field{Last seen} column on the readers tab. \\
Weekly  & Skim the audit trail for variance closures and manual adjustments. \\
Weekly  & Take a backup (see below). \\
Monthly & Review user accounts and disable anyone who has left. \\
\bottomrule
\end{tabular}
\end{table}

\subsection{Backups}

The database is a single file, but it is in use while the server runs, so
\textbf{do not simply copy it}. Use the safe backup command:

\begin{plainbox}
\typed{sqlite3 data/denim\_rfid.db ".backup data/backup-2026-08-18.db"}
\end{plainbox}

Keep the copies somewhere off the machine.

\subsection{Housekeeping the event history}

At full production the tracking history grows by roughly a million rows and
400\,MB a day. The dashboards stay fast regardless, because they read a
pre-summarised table. Decide how long you must keep the raw scan history for
traceability, and archive anything older.
